% TSYP14 Challenges Selection (compile with lualatex)
\documentclass[11pt,a4paper]{article}

\usepackage[top=2.4cm,bottom=2.4cm,left=2.6cm,right=2.6cm]{geometry}
\usepackage{fontspec}
\setmainfont{Lato}[BoldFont={Lato Black}]
\usepackage{xcolor}
\usepackage{enumitem}
\usepackage[hidelinks]{hyperref}
\usepackage{ulem}

\pagestyle{empty}
\setlength{\parindent}{0pt}
\setlength{\parskip}{8pt}
\linespread{1.12}

\definecolor{titleblue}{HTML}{073762}
\definecolor{linkblue}{HTML}{1154CC}

% "More details" links: specification books from tsyp.ieee.tn/program/challenges
\newcommand{\linkAOTN}{https://online.fliphtml5.com/cayrn/ctjg/}
\newcommand{\linkNEXUS}{https://drive.google.com/file/d/1IGdD5hincq2haXRrSGjGAKNVMgv2DQnd/view?usp=sharing}
\newcommand{\linkCITY}{https://drive.google.com/file/d/18_xa-QnqOzZSWuOMNb1TwdR93buSmO_t/view?usp=sharing}
\newcommand{\linkIAS}{https://drive.google.com/file/d/1CZ472UIetrNzyy4dZiYGaYu8XxrVkXQp/view?usp=sharing}
\newcommand{\linkWIE}{https://drive.google.com/file/d/1CE_09tClW84i7t-oh6olvyofaJQwmkJH/view?usp=sharing}
\newcommand{\linkCIS}{https://drive.google.com/file/d/1pocl9K2ymSAoT2N2oNVClJ9B-ER0pKbc/view?usp=sharing}
\newcommand{\linkENERGY}{https://drive.google.com/file/d/15NEu59huEA7_WRwFq4H2khQCtqOV08EY/view?usp=sharing}
\newcommand{\linkCAS}{https://drive.google.com/file/d/1zEz4RCmTzVbJ711nImh2l9uV_QVDBElu/view?usp=sharing}

% \challenge{number}{name}{dark colour}{light colour}
\newcommand{\challenge}[4]{%
  \colorlet{chdark}{#3}\colorlet{chlight}{#4}%
  \vspace{14pt}
  {\leftskip=0.6cm\raggedright\fontsize{15}{18}\selectfont\bfseries\color{chdark}#1.\enspace #2\par}
  \vspace{2pt}}
\newcommand{\sub}[2]{%
  {\hspace*{2cm}\fontsize{13}{16}\selectfont\bfseries\color{chlight}#1.\enspace #2\par}}
\newcommand{\moredetails}[1]{%
  \vspace{4pt}{\hspace*{0.6cm}\href{#1}{\color{linkblue}\uline{More details}}\par}}

\setlist[itemize]{label=\textbullet,leftmargin=1.1cm,itemsep=3pt,topsep=2pt}

\begin{document}

\begin{center}
  {\fontsize{24}{30}\selectfont\bfseries\color{titleblue}TSYP14 CHALLENGES SELECTION\par}
\end{center}
\vspace{6pt}

%% ---------------------------------------------------------------- 1
\challenge{1}{AOTN Challenge: Architects of the Noosphere}{green!45!black}{green!55!black}

\sub{a}{Challenge Description:}

The term \emph{noosphere} describes a layer of collective human thought and reasoning: the way
humanity's shared ideas shape the world. Recent crises, from natural disasters to public health
emergencies, show that technology alone rarely saves a system under stress. What usually breaks
first is the decision-making around it: fragmented information, eroded trust, and shortcuts taken
under pressure.

Architects of the Noosphere challenges student teams to design ethical, collaborative
intelligence systems that strengthen how communities make decisions and recover under pressure.
It is a progressive innovation journey grounded in the UN Sustainable Development Goals, around
three pillars: Sustainable Education, Sustainable Industry and Sustainable Agriculture. Teams
identify a genuine problem, design a solution, build a working prototype, and defend a
\textbf{complete business case} in front of a jury, culminating in a final pitch at the TSYP
Congress. Along the way, unannounced twists test how well teams adapt when the rules change.

\sub{b}{Proposed Solution:}

A \textbf{remote alerting system} that turns hidden warnings from factory transformer substations
into early, actionable alerts, before a slow, detectable fault grows into a fire.

\sub{c}{Domains explored:}
\begin{itemize}
  \item \textbf{Entrepreneurship and Marketing:} To build a sustainable business model, size the
        market and plan how the solution reaches factories.
  \item \textbf{Web Development:} To build the platform that delivers alerts and lets the
        responsible staff follow and confirm them.
  \item \textbf{IoT and Embedded Systems:} To connect the substation's sensors and send their
        warnings out in real time.
  \item \textbf{Soft Skills and Pitching:} To validate the problem with real users and defend the
        project in front of the jury.
\end{itemize}
\moredetails{\linkAOTN}

\newpage
%% ---------------------------------------------------------------- 2
\challenge{2}{NEXUS Challenge: Next-Generation Emergency eXchange for Universal Sustainability}{red!70!black}{red!80!black}

\sub{a}{Challenge Description:}

Tunisia faces increasing climate and infrastructure risks, including flash floods, wildfires,
drought and dam shortfalls, coastal erosion and summer heatwaves. Disaster losses are often not
caused by a lack of data, but by data that does not reach an early-warning system, responder or
household in time to act.

The challenge focuses on building low-cost, low-power smart systems that connect sensing,
communication and response into one resilient system: \textbf{Detect, Communicate, Coordinate,
Enable Response}. Solutions must be designed for Tunisia's reality, with intermittent power,
limited connectivity, local maintenance and affordability, and must get information to the
household, municipality, Civil Protection or farmer in a usable form and in time to act. Teams are
encouraged to address a real Tunisian region and to connect their solution to the UN Sustainable
Development Goals.

\sub{b}{Proposed Solution:}

A \textbf{low-cost, offline-capable flash-flood early warning chain} for Tunisian wadi crossings:
it detects a rising wadi, delivers the warning through whichever channel still works, and
escalates automatically when nobody confirms it.

\sub{c}{Domains explored:}
\begin{itemize}
  \item \textbf{IoT and LoRaWAN:} To connect solar-powered sensor nodes along the wadi over
        long-range, low-power radio, without depending on the mobile network.
  \item \textbf{Embedded Systems:} To build the nodes that measure water level and trigger the
        local beacon and siren on their own.
  \item \textbf{Entrepreneurship:} To build the business and humanitarian case: costs, funding,
        stakeholders and community impact.
\end{itemize}
\moredetails{\linkNEXUS}

\newpage
%% ---------------------------------------------------------------- 3
\challenge{3}{City Brain Convergence Challenge}{blue!60!black}{blue!70!black}

\sub{a}{Challenge Description:}

Smart cities depend on sensing, deciding and trusting all working together, yet most systems are
built to do only one well. This challenge splits a smart-city system into three layers: the
\textbf{Twin} (which simulates the city's sensors and actuators), the \textbf{City Brain} (which
consumes real-time data and makes operational decisions) and the \textbf{Guardian} (which monitors
the same data to detect anomalies, malfunctions and cyber-attacks).

Teams are grouped into triplets, one team per layer, and each team builds its layer independently
under a shared technical contract. The three systems must then interoperate as one functioning
smart city and be tested under realistic and adversarial conditions, including when the data itself
is being falsified. The challenge emphasizes engineering to an interface you do not control, and
proving that your system remains reliable when connected to systems you did not build.

\sub{b}{Proposed Solution:}

The \textbf{decision-making layer of a smart city}: it turns live sensor data into incidents and
chooses safe actions using a local LLM, weighted by security trust scores.

\sub{c}{Domains explored:}
\begin{itemize}
  \item \textbf{Agentic AI and LLMs:} To reason about each incident and propose the right response
        with a language model running locally on an edge device.
  \item \textbf{Graph Databases:} To store the city map and link related events and incidents
        together, giving the AI the context it needs.
  \item \textbf{IoT and MQTT:} To exchange secure, real-time messages with the Twin and Guardian
        layers built by the other teams.
\end{itemize}
\moredetails{\linkCITY}

\newpage
%% ---------------------------------------------------------------- 4
\challenge{4}{IAS Challenge: Engineering Industry for a Dynamic Future}{orange!80!black}{orange!90!black}

\sub{a}{Challenge Description:}

Industrial systems are designed to operate efficiently under defined conditions. However,
real-world industrial environments continuously evolve due to changing demands, unexpected
disruptions, resource constraints, operational challenges and technological advances. The IAS
Challenge focuses on developing innovative engineering solutions that address real industrial
needs, with emphasis on \textbf{resilience, efficiency, sustainability and digitalisation}.

Teams start from a real industrial problem or process and follow a research- and process-based
methodology: industrial problem, process mapping, research and state of the art, solution
exploration, comparative analysis, solution development, validation and impact assessment. The
goal is to deliver practical, technically sound solutions that move an industrial process from its
current AS-IS state to an improved TO-BE state, and prepare industrial systems for a more dynamic
and demanding future.

\sub{b}{Proposed Solution:}

An \textbf{AI-driven IoT system} for smarter, safer and more productive fish hatcheries.

\sub{c}{Domains explored:}
\begin{itemize}
  \item \textbf{IoT and Embedded Systems:} To monitor water conditions and hatchery equipment
        continuously in real time.
  \item \textbf{AI and ML:} To predict risks and optimize the hatchery process from the collected
        data.
  \item \textbf{Computer Vision:} To observe the fish and detect abnormal behaviour or conditions
        automatically.
  \item \textbf{Automation and Control:} To act on the process automatically and keep conditions
        within safe limits.
\end{itemize}
\moredetails{\linkIAS}

\newpage
%% ---------------------------------------------------------------- 5
\challenge{5}{WIE Challenge: Women and Rural Agriculture Innovation}{magenta!60!black}{magenta!70!black}

\sub{a}{Challenge Description:}

This challenge invites engineering and entrepreneurship students to design technology-driven and
business-driven solutions that strengthen the position of women in Tunisia's rural and peri-urban
agricultural and artisanal economy. Women agricultural workers, artisanal producers and
micro-entrepreneurs play a major role in local production, such as olive oil harvesting, essential
oil distillation, textile weaving and agro-food processing.

However, they face interconnected challenges: low technological automation, fragmented and
informal supply chains, occupational health risks, limited direct access to markets, and dependence
on intermediaries. Teams choose one track (reducing physical drudgery, structuring supply chains and
market access, or improving occupational safety) and develop practical, deployable solutions with
a realistic business concept for rural Tunisian communities.

\sub{b}{Proposed Solution:}

A \textbf{digital market-access and supply-chain platform} connecting rural women farmers
directly to customers, with traceability, fair price transparency and offline access.

\sub{c}{Domains explored:}
\begin{itemize}
  \item \textbf{Web and Mobile Development:} To build an offline-first platform that also works
        on basic phones through SMS, and handles orders, collection and traceability.
  \item \textbf{UI/UX:} To design simple, accessible interfaces adapted to producers with limited
        digital experience.
  \item \textbf{Entrepreneurship:} To build a viable business model with fair income for the
        producers and realistic partnerships.
\end{itemize}
\moredetails{\linkWIE}

\newpage
%% ---------------------------------------------------------------- 6
\challenge{6}{Ground Truth Under Attack: Securing the Earth Observation Chain}{teal!70!black}{teal!80!black}

\sub{a}{Challenge Description:}

Earth-observation products support decisions in agriculture, insurance, water management and
disaster response. Their trustworthiness depends on evidence of origin, time, integrity and
processing history across the sensing, communication and product chain. Yet documented navigation
interference, false-data injection and increasingly capable image synthesis show that
plausible-looking data may still be untrustworthy, and a failure may remain silent while a
corrupted observation reaches the analyst.

Hosted jointly by the IEEE CIS, GRSS and SPS chapters, the challenge asks teams to select one
segment (ground link, in-situ sensing or imagery) and build a verifiable defensive prototype focused
on \textbf{provenance, integrity verification, anomaly detection and controlled adversarial
testing}, using authentic public data and reporting detection quality, false alarms, latency and
security overhead.

\sub{b}{Proposed Solution:}

An \textbf{AI-driven security system} detecting false data injected into Earth observation
sensors.

\sub{c}{Domains explored:}
\begin{itemize}
  \item \textbf{Cybersecurity:} To sign data at the source and prove it has not been tampered with
        along the chain.
  \item \textbf{AI and ML:} To detect falsified readings that look plausible, and test the system
        against a simulated attacker.
  \item \textbf{Web Development:} To build the dashboard that shows alerts and verification
        results to the analyst.
\end{itemize}
\moredetails{\linkCIS}

\newpage
%% ---------------------------------------------------------------- 7
\challenge{7}{EnergyMind Challenge: Architects of the Energy Noosphere}{olive!70!black}{olive!80!black}

\sub{a}{Challenge Description:}

Tunisia's grid faces growing strain from seasonal peaks, fast-expanding renewable energy, and
increasing numbers of smart meters and IoT devices. These create more points where a fault,
miscalibration or biased reading can silently corrupt the data used by automated systems. As AI
takes on more grid decisions, the question is not only whether it predicts well, but also: can it be
trusted, and who is accountable when it is wrong?

The EnergyMind Tunisia challenge invites participants to build the \textbf{trust layer} of
Tunisia's energy systems: an AI that stands between raw energy data and automated action. It must
detect when data or equipment is failing, faulty or falsified, distinguish genuine grid events from
data-quality problems, explain every verdict with a confidence score, decide transparently when it
can act alone and when a human must decide, and design how multiple nodes could share what they
learn while protecting privacy.

\sub{b}{Proposed Solution:}

An \textbf{explainable AI system} detecting anomalies in the smart grid and escalating critical
decisions to humans.

\sub{c}{Domains explored:}
\begin{itemize}
  \item \textbf{AI and ML:} To detect anomalies in grid data, give each verdict a confidence score
        and explain it in plain language.
  \item \textbf{UI/UX:} To design the operator interface where humans review, approve or override
        the AI's decisions quickly.
\end{itemize}
\moredetails{\linkENERGY}

\newpage
%% ---------------------------------------------------------------- 8
\challenge{8}{CAS Challenge: Edge AI \& IoT for Resilient Smart Infrastructure}{violet!70!black}{violet!80!black}

\sub{a}{Challenge Description:}

The CAS Challenge focuses on Edge AI, IoT, embedded systems, smart sensing, wireless communication,
RFID and intelligent systems. Teams are invited to develop solutions for resilient infrastructure and
real-world applications by combining sensing, intelligent processing, connectivity and real-time
action.

Teams must develop an intelligent connected system that addresses a real-world problem by
\textbf{collecting information, processing it intelligently at the edge, and producing an
appropriate response or action}. A meaningful processing capability must run on the device itself:
cloud-only solutions are not sufficient. The final solution is demonstrated live through a
functional prototype showing sensing, intelligent processing, decision-making and action.

\sub{b}{Proposed Solution:}

An \textbf{edge AI device} translating sign language into text and voice in real time.

\sub{c}{Domains explored:}
\begin{itemize}
  \item \textbf{Computer Vision:} To recognize hand gestures and signs from a camera feed.
  \item \textbf{Edge AI:} To run the recognition model directly on the device in real time,
        without the cloud.
  \item \textbf{Embedded Systems:} To build a compact, low-cost device that captures, translates
        and speaks.
\end{itemize}
\moredetails{\linkCAS}

\newpage
%% ---------------------------------------------------------------- 9
\challenge{9}{Booth Project Challenge}{gray!60!black}{gray!70!black}

The Booth Project Challenge gives us the opportunity to present our work to professionals and peers
from other schools. It's a chance to showcase our creativity, problem-solving skills, and the impact
of our project, while receiving valuable feedback and learning from the approaches of others.

\sub{a}{Domains Explored:}
\begin{itemize}
  \item \textbf{Embedded systems:} To build hardware-integrated solutions that perform intelligent
        sensing or automation tasks.
  \item \textbf{IoT:} To connect smart devices for data collection, analysis and real-time control.
  \item \textbf{Video game dev:} To engage users or educate audiences through interactive
        experiences that blend creativity with technology.
  \item \textbf{etc.:} Open to innovative technologies like robotics, AR/VR\ldots
\end{itemize}

\end{document}
